\section{Discrete Valuation Rings and Dedekind Domains}

\subsection{Exercise - Complete}\label{9.1}

\begin{proof}
Since $A$ is a Dedekind domain, it is a Noetherian domain of dimension 1. Thus, $S^{-1}A$ is a Noetherian domain (Proposition 7.3).
Since all the prime ideals of $S^{-1}A$ are extended prime ideals of $A$ (Proposition 3.11), a chain of prime ideals 
$S^{-1}\mathfrak{p}_0 \subset S^{-1}\mathfrak{p}_1 \subset S^{-1}\mathfrak{p}_2$ pulls back to a chain of prime ideals
$\mathfrak{p}_0 \subset \mathfrak{p}_1 \subset \mathfrak{p}_2$ of $A$. Then, since $A$ has dimension 1, $S^{-1}A$ has dimension at
most 1.

Suppose $S^{-1}A$ has dimension $0$. Since $S^{-1}A$ is an integral domain, $0$ is a prime. Every $\mathfrak{p} \in \Spec(S^{-1}A)$
satisfies $0 \subset \mathfrak{p}$, so since $S^{-1}A$ has dimension $0$, $\mathfrak{p} = 0$. Thus, the only primes of $S^{-1}A$
are $0$, so $S^{-1}A$ is a field. In particular, every prime ideal of $A$ intersects $S$. Let $\mathfrak{a}$ be an arbitrary non-zero
ideal of $A$. Since $A$ is a domain, $\mathfrak{a}^e \neq 0$, so $\mathfrak{a}^e = (1)$, implying $\mathfrak{a}$ intersects $S$ by
Proposition 3.11. Now, let $a \in A$ and let $b \in A$ be such that $b \neq 0$; this yields an arbitrary element $a/b \in K$. Since $(b)$ 
is nonzero, $(b)$ intersects $S$. Thus, there is some $x \in A$ such that $xb \in S$. Thus, $\frac{a}{b} = \frac{xa}{xb} \in S^{-1}A$.
Since $A \subset S^{-1}A \subset K$, we see that $S^{-1}A = K$.

Alternatively, $S^{-1}A$ has dimension 1, so it is a Noetherian domain of dimension 1. Since $A$ is integrally closed in its field of fractions $K$,
and $A \subset S^{-1}A \subset K$, we have by Proposition 5.12 that $S^{-1}A$ is closed in its field of fractions, also $K$. Thus, $S^{-1}A$
is integrally closed, so by the equivalent conditions of Theorem 9.3, $S^{-1}A$ is a Dedekind domain.
\end{proof}

\begin{proof}
Let $H$ and $H'$ be the ideal class groups of $A$ and $S^{-1}A$, respectively. If $S^{-1}A$ is a field, then $H' = 1$, so there is clearly a surjective
homomorphism $H \to H'$. Otherwise, by the previous proof, $S^{-1}A$ is a Dedekind domain. 

Note that $I'$, the group of fractional ideals of $S^{-1}A$, is generated by the prime ideals of $S^{-1}A$. The ring homomorphism $A \to S^{-1}A$ naturally 
gives rise to a group homomorphism $\mathfrak{a} \to \mathfrak{a}^e$. By Proposition 3.11, all prime ideals of $S^{-1}A$ are extended, so since the prime
ideals generate $I'$, the homomorphism $I \to I'$ is surjective.

Similarly, if $xA \in P$, then $xA \mapsto (x/1)S^{-1}A \in P'$, so the map $I \to I'$ restricted to $P$ gives a homomorphism $P \to P'$. Thus, our homomorphism
lifts to $H = I/P \to I'/P' = H'$, and it is surjective since $I \to I'$ is surjective.
\end{proof}

\subsection{Exercise - Complete}\label{9.2}

\begin{proof}
For a polynomial $f = \sum_i a_i x^i$ with coefficients in a Dedekind domain, we define the ``content'' of $f$ to be the ideal $c(f) = (a_i)_i$ of $A$, the ideal generated 
by its coefficients. Now, suppose $f,g \in A[x]$, where $f = \sum a_i x^i$ and $g = \sum b_j x^j$. If $c(fg) = 0$, then $fg = 0$ so that either $f = 0$ or $g = 0$; $A$ is an
integral domain, so $A[x]$ is also. Then, 
$c(f) = 0$ or $c(g) = 0$ so that $c(fg) = c(f)c(g)$. If $c(f)c(g) = 0$, then by containment, $c(fg) = 0$. Thus, we can assume that $c(fg)$, $c(f)$, and $c(g)$ are 
nonzero ideals. Then, since $A$ is a Dedekind domain, each is uniquely described as a product of nonzero prime ideals. In particular, to show that $c(fg) = c(f)c(g)$, we must
show that for each nonzero prime ideal $\mathfrak{p}$, the exponent of $\mathfrak{p}$ of $c(fg)$ is the sum of its exponents in $c(f)$ and $c(g)$. 

Let's localize at an arbitrary nonzero prime ideal $\mathfrak{p}$ of $A$. Since $A$ is Dedekind, every nonzero prime ideal is maximal. In particular, any nonzero prime that is not
$\mathfrak{p}$ extends to all of $A_{\mathfrak{p}}$. Thus, under $A \mapsto A_{\mathfrak{p}}$, the factorizations of $c(fg)$, $c(f)$, and $c(g)$ map to
$(S^{-1}\mathfrak{p})^{v(fg)}$, $(S^{-1}\mathfrak{p})^{v(f)}$, and $(S^{-1}\mathfrak{p})^{v(g)}$, respectively. These exponents are the same as the exponents on $\mathfrak{p}$
in the original factorizations. Thus, we must show that $v(fg) = v(f) + v(g)$. By Theorem 9.3, $A_{\mathfrak{p}}$ is a discrete valuation ring since $A$ is a Dedekind ring.
Thus, we've reduced the problem to the following: Let $f = \sum_i a_i x^i$, $g = \sum_j b_j x^j \in A[x]$ be nonzero, where $A$ is a discrete valuation ring. Then, we must show 
that $c(fg) = c(f)c(g)$.

In this new setting, since $A$ is a discrete valuation ring, there is some $\pi \in A$ such that every nonzero ideal has the form $(\pi)^k$, $k \ge 0$. Thus, $c(f) = (\pi)^{v(f)}$ 
and $c(g) = (\pi)^{v(g)}$. Since $c(f) = (a_i)_i$ and $c(g) = (b_j)_j$, for all $a_i$ and $b_j$, $a_i \in (\pi)^{v(f)}$ and $b_j \in (\pi)^{v(g)}$. Thus, 
$a_i' = a_i/\pi^{v(f)} \in A$ and $b_j' = b_j/\pi^{v(g)} \in A$. Let $f_0 = \sum a_i' x^i$ and $g_0 = \sum b_j' x^j$. Then, $f_0 = f / \pi^{v(f)}$ and $g_0 = g / \pi^{v(g)}$ so that 
$c(f_0) = (a_i')_i = (a_i/\pi^{v(f)})_i = (a_i)_i/(\pi)^{v(f)} = (\pi)^{v(f)}/(\pi)^{v(f)} = (1)$; similarly, $c(g_0) = (1)$. By Exercise \ref{1.2}, $c(f_0g_0) = (1)$. Finally, 
$fg = \pi^{v(f)} f_0 \pi^{v(g)} g_0 = \pi^{v(f) + v(g)} f_0 g_0$. Then, 
$$c(fg) = (\pi)^{v(f) + v(g)}c(f_0g_0) = (\pi)^{v(f) + v(g)} = (\pi)^{v(f)}(\pi)^{v(g)} = c(f)c(g),$$
as desired.
\end{proof}

\subsection{Exercise - Complete}\label{9.3}

\begin{proof}
If $A$ is a discrete valuation ring, then $A$ is Noetherian by definition. Conversely, suppose $A$ is a Noetherian valuation ring that is not a field. 
Then, $A$ is an integral domain, by definition, and by Proposition 5.18, $A$ is a local ring that is integrally closed. Supposing for a moment that $A$
has dimension one, then $A$ is a local Noetherian domain of dimension one that is integrally closed, so by Proposition 9.2, $A$ is a discrete valuation ring.
Therefore, we just have to show that all nonzero prime ideals of $A$ are maximal.

Since $A$ is a local ring, it has a unique maximal ideal $\mathfrak{m}$. Let $\mathfrak{p}$ be an arbitrary prime ideal of $A$ so $\mathfrak{p} \subset \mathfrak{m}$. 
Using the construction of Exercise \ref{5.30}, we have a valuation $v : K^{\times} \to K^{\times}/U$, where $K$ is the field of fractions of $A$ and $U$ is the group of 
units of $A$.

Now, we will construct an ascending chain of ideals. Let $y \in \mathfrak{p}$ and $\mathfrak{a}_1 = \{a \in A \mid v(a) \ge v(y)\}$.
If $a \in A$ is such that $v(a) \ge v(y)$, then $ay^{-1} \in A$, so $a = ay^{-1} \cdot y \in \mathfrak{p}$. Thus, $\mathfrak{a}_1 \subset \mathfrak{p}$. Suppose there
is some $y' \in \mathfrak{p} - \mathfrak{a}_1$. We define $\mathfrak{a}_2 = \{a \in A \mid v(a) \ge v(y')\}$. We have a proper inclusion $\mathfrak{a}_1 \subset \mathfrak{a}_2$ 
since $y' \in \mathfrak{a}_2 - \mathfrak{a}_1$, implying $v(y') < v(y)$. By the same argument as before, $\mathfrak{a}_2 \subset \mathfrak{p}$. Continue in this fashion; 
since $A$ is Noetherian, there is some $n$ such that the strictly ascending chain of ideals terminates. Then, $\mathfrak{p} - \mathfrak{a}_n$ is empty so that 
$\mathfrak{p} = \mathfrak{a}_n$. Thus, there is some $y \in A$ such that $\mathfrak{p} = \{a \in A \mid v(a) \ge v(y)\}$.

Suppose to the contrary that $\mathfrak{p}$ is not $\mathfrak{m}$, so there is some $x \in \mathfrak{m} - \mathfrak{p}$, which is clearly nonzero. Since $x \in \mathfrak{m}$, 
$v(x) > 0$. Because $x \not\in \mathfrak{p}$ and $y \in \mathfrak{p}$, $x \not\in (y)$ so $xy^{-1} \not\in A$. Since $A$ is a valuation ring, we must then have $yx^{-1} \in A$.
Further, $v(yx^{-1}) = v(y) - v(x) < v(y)$ since $v(x) > 0$ so $yx^{-1} \not\in \mathfrak{p}$. Since $x, yx^{-1} \in A-\mathfrak{p}$ and $\mathfrak{p}$ is prime, we have
$y = x \cdot yx^{-1} \not\in \mathfrak{p}$. This is a clear contradiction. Thus, the only nonzero prime ideal of $A$ is the unique maximal ideal, so all nonzero prime ideals
of $A$ are maximal, as desired.
\end{proof}

\subsection{Exercise - Complete}\label{9.4}

\begin{proof}
Let $A$ be a local domain (not a field) in which the maximal ideal $\mathfrak{m}$ is principal and $\bigcap_{i=1}^{\infty}\mathfrak{m}^n = 0$. Suppose for a moment
that $\mathfrak{m}$ is the only nonzero prime ideal. Then, every prime ideal is finitely generated since $\mathfrak{m}$ is principal, so by Exercise \ref{7.1}, $A$ 
is Noetherian. Further, it has dimension one since $0 \subset \mathfrak{m}$ is the only strict chain of primes. Then, since $A$ is a Noetherian local domain of 
dimension one with a maximal ideal that is principal, we have by Proposition 9.2 that $A$ is a discrete valuation ring.

Therefore, we just have to show that any nonzero prime ideal $\mathfrak{p}$ of $A$ is equal to $\mathfrak{m}$. Since $\mathfrak{m}$ is principal, $\mathfrak{m} = (\pi)$ 
for some $\pi \in A$. Now, let $y \in \mathfrak{p}$. If $y = 0$, then $y \in \mathfrak{m}$, so suppose $y \neq 0$. By assumption, $\bigcap_{i=1}^{\infty}\mathfrak{m}^n = 0$, 
so $y \not\in \bigcap_{i=1}^{\infty}\mathfrak{m}^n$. Thus, there is some minimal $n(y) > 1$ (since $y \in \mathfrak{m}$) such that $y \not\in \mathfrak{m}^{n(y)}$. Then, 
$\{n(y) \mid y \in \mathfrak{p} - 0\}$ is a non-empty subset of integers greater than $1$, so it has a minimum $N > 1$. Thus, there is some $b \in \mathfrak{p}$ such that 
$b \not\in \mathfrak{m}^N$. Since $b \in \mathfrak{p} \subset \mathfrak{m} = (\pi)$, we can write $b = a \pi$, where $a \in A$. 

If $a \in \mathfrak{m}^{N-1} = (\pi^{N-1})$, then $b = a \pi \in (\pi^N)$, contradicting our choice of $b$, so $a \not\in \mathfrak{m}^{N-1}$.
Since $b \in \mathfrak{p}$, either $\pi \in \mathfrak{p}$ or $a \in \mathfrak{p}$. But $a \not\in \mathfrak{p}$; otherwise, $a \in \mathfrak{p}$ 
and $a \not\in \mathfrak{m}^{N-1}$ together violate the minimality of $N$. Thus, $\pi \in \mathfrak{p}$ so that $\mathfrak{m} = (\pi) \subset \mathfrak{p}$, 
implying that $\mathfrak{p} = \mathfrak{m}$, as desired.
\end{proof}

\subsection{Exercise - Complete}\label{9.5}

\begin{proof}
Let $M$ be a finitely generated module over a Dedekind domain. Since $M$ is finitely generated and $A$ is Noetherian, we have by Exercise \ref{7.16} that
$M$ is flat if and only if $M_{\mathfrak{m}}$ is free for all maximal ideals $\mathfrak{m}$ of $A$. By Exercise \ref{3.13}, $M$ is torsion-free if and only 
if $M_{\mathfrak{m}}$ is torsion-free for all maximal ideals $\mathfrak{m}$ of $A$. Thus, to show that $M$ is flat if and only if $M$ is torsion-free, it
suffices to show that for all maximal ideals $\mathfrak{m}$ of $A$, $M_{\mathfrak{m}}$ is torsion-free if and only if $M_{\mathfrak{m}}$ is free.

Note that if $M_{\mathfrak{m}} = 0$, the result is trivial, so suppose $M_{\mathfrak{m}} \neq 0$. Since $M_{\mathfrak{m}}$ is finitely generated,
if $M_{\mathfrak{m}}$ is free, $M_{\mathfrak{m}} \cong A_{\mathfrak{m}}^n$ as an $A_{\mathfrak{m}}$-module for some $n > 0$. Let 
$(a_1,\ldots,a_n) \in A_{\mathfrak{m}}^n$, and let $x \in A_{\mathfrak{m}}$ be non-zero. If $x(a_1,\ldots,a_n) = 0$, then $(xa_1,\ldots,xa_n) = (0,\ldots,0)$.
Since $A_{\mathfrak{m}}$ is an integral domain and $x \neq 0$, $a_i = 0$ for all $i$, so $(a_1,\ldots,a_n) = (0,\ldots,0)$. Thus, $T(M_{\mathfrak{m}}) = 0$.

Conversely, suppose $T(M_{\mathfrak{m}}) = 0$. Since $A$ is a Dedekind domain, we have by Theorem 9.3 that $A_{\mathfrak{m}}$ is a discrete valuation ring. Thus, 
by Proposition 9.2, $A_{\mathfrak{m}}$ is a principal ideal domain. Hence, by the structure theorem for finitely generated modules over a principal ideal domain,
$$M_{\mathfrak{m}} \cong A_{\mathfrak{m}}^n \oplus A_{\mathfrak{m}}/\mathfrak{a}_1 \oplus \cdots \oplus A_{\mathfrak{m}}/\mathfrak{a}_r$$
for nonzero ideals $\mathfrak{a}_i$ of $A_{\mathfrak{m}}$, where $r \ge 0$. If $r \ge 1$, $M_{\mathfrak{m}}$ has an $A_{\mathfrak{m}}$-submodule $M'$ isomorphic to
$A_{\mathfrak{m}}/\mathfrak{a}_1$. Let $m \in M'$ be the preimage of $1$ in the isomorphism $M' \cong A_{\mathfrak{m}}/\mathfrak{a}_1$, and let $a \in \mathfrak{a}_1$
be a nonzero element. Then, $a m = 0$, violating the fact that $M_{\mathfrak{m}}$ is torsion-free. Thus, $r = 0$ so that $M_{\mathfrak{m}} \cong A_{\mathfrak{m}}^n$,
as desired.
\end{proof}

\subsection{Exercise - Complete}\label{9.6}

\begin{proof}
Let $\mathfrak{p}$ be a nonzero prime ideal of $A$. Since $A$ is a Dedekind domain, $A_{\mathfrak{p}}$ is a discrete valuation ring by Theorem 9.3. Thus, $A_{\mathfrak{p}}$
is a principal ideal domain by Proposition 9.2. Suppose $M_{\mathfrak{p}} \neq 0$. By the structure theorem for principal ideal domains,
$$M_{\mathfrak{p}} \cong A_{\mathfrak{p}}^n \oplus A_{\mathfrak{p}}/\mathfrak{a}_1 \oplus \cdots \oplus A_{\mathfrak{p}}/\mathfrak{a}_r,$$
where $\mathfrak{a}_i$ are nontrivial ideals of $A_{\mathfrak{p}}$. By Exercise \ref{3.13}, $T(M_{\mathfrak{p}}) = T(M)_{\mathfrak{p}} = M_{\mathfrak{p}}$. If $n \ge 1$, there
is a submodule $M' \cong A_{\mathfrak{p}}$ of $M_{\mathfrak{p}}$. Since $A_{\mathfrak{p}}$ is an integral domain, $M'$ contains non-torsion elements, contradicting $T(M_{\mathfrak{p}}) = M_{\mathfrak{p}}$.
Thus, $n = 0$ so that
$$M_{\mathfrak{p}} \cong A_{\mathfrak{p}}/\mathfrak{a}_1 \oplus \cdots \oplus A_{\mathfrak{p}}/\mathfrak{a}_r.$$
Since $A_{\mathfrak{p}}$ is a discrete valuation ring with maximal ideal $\mathfrak{p}$, we have $\mathfrak{a}_j = \mathfrak{p}^{n_j}$ for some $n_j > 0$ (Proposition 9.2). Thus,
$$M_{\mathfrak{p}} \cong A_{\mathfrak{p}}/\mathfrak{p}^{n_1} \oplus \cdots \oplus A_{\mathfrak{p}}/\mathfrak{p}^{n_r}.$$
Note that this decomposition is unique (modulo reordering, of course) by the structure theorem. Let $n(\mathfrak{p}) = \max\{n_1,\ldots,n_r\}$. Then clearly, 
$\Ann(M_{\mathfrak{p}}) = \mathfrak{p}^{n(\mathfrak{p})}$. In particular, note that $\Ann(M_{\mathfrak{p}}) \neq 0$ whenever $M_{\mathfrak{p}} \neq 0$.

Since $M$ is finitely generated, by Proposition 3.14, $\Ann(M)_{\mathfrak{p}} = \Ann(M_{\mathfrak{p}}) = \mathfrak{p}^{n(\mathfrak{p})}$. If $\Ann(M) = 0$, then
$\Ann(M_{\mathfrak{p}}) = 0$ so that $M_{\mathfrak{p}} = 0$ for all prime ideals $\mathfrak{p}$. Thus, by Proposition 3.8, $M = 0$. So, if $M \neq 0$, $\Ann(M)$
is a nonzero ideal. Since $\Ann(M)$ is a nonzero ideal of a Dedekind domain $A$, $\Ann(M)$ has a unique decomposition as a product of prime powers (Corollary 9.4). If
$\mathfrak{p}$ is a nonzero prime that is not in the decomposition, then since $A$ is Dedekind, $\mathfrak{p}$ does not contain any of the primes in the decomposition since
all nonzero primes all maximal. Thus, $0 = \Ann(M)_{\mathfrak{p}} = \Ann(M_{\mathfrak{p}})$, so $M_{\mathfrak{p}} = 0$. Thus, $\Supp(M)$ is finite.

Let $\{\mathfrak{p}_1,\ldots,\mathfrak{p}_k\} = \Supp(M)$. For each $\mathfrak{p}_i$, we have, as above,
$$M_{\mathfrak{p}_i} \cong A_{\mathfrak{p}_i}/\mathfrak{p}_i^{n_{i,1}} \oplus \cdots \oplus A_{\mathfrak{p}_i}/\mathfrak{p}_i^{n_{i,r_i}}.$$
Note that we have an $A$-module homomorphism $\phi : M \to \bigoplus_{i=1}^k \bigoplus_{j=1}^{r_i} A/\mathfrak{p}_i^{n_{i,j}}$. If $\mathfrak{p} \not\in \Supp(M)$,
then we have $\phi_{\mathfrak{p}} : M_{\mathfrak{p}} \to \bigoplus_{i=1}^k \bigoplus_{j=1}^{r_i} A_{\mathfrak{p}}/\mathfrak{p}_i^{n_{i,j}}$. For all $\mathfrak{p}_i \in \Supp(M)$,
$\mathfrak{p}$ does not contain $\mathfrak{p}_i$, so $A_{\mathfrak{p}}\mathfrak{p}_i^{n_{i,r_i}} = A_{\mathfrak{p}}$ so that all quotients on the right-hand side vanish. Further,
$M_{\mathfrak{p}} = 0$, so $\phi_{\mathfrak{p}} : 0 \to 0$, which is clearly an isomorphism. On the other hand, let $\mathfrak{p}_s \in \Supp(M)$. Then,
$\phi_{\mathfrak{p}_s} : M_{\mathfrak{p}_s} \to \bigoplus_{i=1}^k \bigoplus_{j=1}^{r_i} A_{\mathfrak{p}_s}/\mathfrak{p}_i^{n_{i,j}}$. If $i \neq s$, then $\mathfrak{p}_s$ does
not contain $\mathfrak{p}_i$, so $A_{\mathfrak{p}_s}\mathfrak{p}_i^{n_{i,j}} = A_{\mathfrak{p}_s}$. In particular, $A_{\mathfrak{p}_s}/\mathfrak{p}_i^{n_{i,j}} = 0$ for $i \neq s$.
Thus, we have $\phi_{\mathfrak{p}_s} : M_{\mathfrak{p}_s} \to \bigoplus_{j=1}^{r_s} A_{\mathfrak{p}_s}/\mathfrak{p}_s^{n_{s,j}}$, which is an isomorphism by construction. By
Proposition 3.9, $\phi$ is an isomorphism, so 
$$M \cong \bigoplus_{i=1}^k \bigoplus_{j=1}^{r_i} A/\mathfrak{p}_i^{n_{i,j}}.$$

To see that this decomposition is unique, note that two decompositions of $M$ must disagree for some prime $\mathfrak{p}$, i.e. the two compositions do not have either the same 
number of factors or the powers of $\mathfrak{p}$ do not match. But then, they would disagree under localization to $M_{\mathfrak{p}}$, which violates the uniqueness guaranteed 
by the structure theorem.
\end{proof}

\subsection{Exercise}\label{9.7}



\subsection{Exercise}\label{9.8}

\subsection{Exercise}\label{9.9}
